%% kapa: exam=1 round=29 year=2018 period=2 form=A subject=accounting subject_ko=회계학 seq=36 number=76
%% source: Q-Net 감정평가사 기출문제 (cst00309) artlSeq=5202808 file=extracted/1008/5202808/2018년 제29회 감정평가사 1차시험 기출문제/2018년 제29회 감정평가사 1차시험 기출문제/2교시 A형.pdf
%% conversion: kordoc 4.24.51 -> tools/gen1.py (auto); review=proofread-sonnet; flags=none
\begin{question}{76}
정상원가계산을 사용하는 (주)감평은 단일제품을 제조ㆍ판매하는 기업이다. 20x1년도의 고정제조간접원가 총예산액 및 실제 발생액은 ￦720,000이었다. 20x1년 제품의 생산 및 판매량은 다음과 같고, 기초 및 기말 재공품은 없다.
\begin{framedtext}
\begin{tabular}{@{}p{0.2\linewidth}p{0.22\linewidth}p{0.16\linewidth}l@{}}
기초재고 & 40,000단위 & 생산량 & 140,000단위 \\
판매량 & 160,000단위 & & \\
\end{tabular}
\end{framedtext}
고정제조간접원가배부율은 120,000단위를 기준으로 산정하며, 이 배부율은 매년 동일하게 적용된다. 한편, 제조원가의 원가차이는 전액 매출원가에서 조정한다. 변동원가계산에 의한 영업이익이 ￦800,000인 경우, 전부원가계산에 의한 영업이익은? \par
\begin{choices}
  \item ￦680,000
  \item ￦700,000
  \item ￦750,000
  \item ￦830,000
  \item ￦920,000
\end{choices}
\end{question}
\answer{1}
